% 03_capi_sampling_quest
\parttitle{Part II: Baseline Household Survey (CAPI)}

The CAPI baseline is the first stage of the L2Phl programme and the foundation for
everything that follows. Conducted face-to-face between September and October 2025,
it produces the nationally representative, household-level estimates that anchor the
project, and it serves as the sampling frame from which the CATI panel draws its
follow-up respondents. The baseline is frozen once fieldwork closes: its statistics
do not change as the panel rolls forward, so it provides a stable reference against
which subsequent rounds are measured.

\section{Sampling}

The baseline uses a two-stage stratified cluster design. The primary sampling unit
(PSU) is the barangay, which serves as the enumeration area; the second-stage unit
is the household. The realized sample covers 2,470 households and 10,496 members
across 247 PSUs.

Strata are defined geographically. Most regions contribute two strata, an urban
stratum and a rural stratum, so the 17 regions yield 34 strata. Five special
geographic areas in the Bangsamoro Autonomous Region in Muslim Mindanao (BARMM),
carrying region codes 101 through 105, each enter as a single stratum. Together
these give 39 strata in total.

Within this frame, PSUs are selected with probability proportional to size (PPS),
and households are then selected within each sampled PSU. This design concentrates
fieldwork in a manageable number of enumeration areas while preserving national
representativeness after weighting.

Three analytical weights accompany the data, each suited to a different level of
analysis. The individual weight scales person-level indicators, so that estimates
such as employment, educational attainment, and health status represent the adult
or member population. The household weight scales household-level behaviour and
decisions, such as finances, asset ownership, and household perceptions. The
population weight expresses results in terms of the number of Filipinos affected,
supporting welfare and SDG-style framing of the kind "X\% of Filipinos live in
households where\ldots". The full construction of these weights, including the base
weight, the exact measure-of-size correction, post-stratification to census
household totals, and age--sex person calibration, is documented in the weighting
technical note.

\begin{refbox}
\rd{CAPI/Analysis/SL/tex/L2PHL_WEIGHTING_TECHNICAL_NOTE.tex}\\
\rd{CAPI/Analysis/SL/tex/L2PHL_WEIGHTING_TECHNICAL_NOTE.pdf}
\end{refbox}

\begin{notebox}{Which weight, in one line}
Use the individual weight for person-level indicators (employment, education,
health); the household weight for household decisions (finance, assets,
perceptions); and the population weight for welfare and SDG framing
("X\% of Filipinos live in households where\ldots").
\end{notebox}

\section{Questionnaire and Modules}

The baseline instrument comprises 14 modules. Each module has a primary unit of
analysis, which determines the weight applied to its indicators.

\begin{table}[h]
\centering
\small
\begin{tabular}{@{}lll@{}}
\toprule
\textbf{Module} & \textbf{Name} & \textbf{Primary unit / weight} \\
\midrule
M00 & Passport   & --- \\
M01 & Roster     & Individual \\
M02 & Education  & Individual \\
M03 & Employment & Individual \\
M04 & Income     & Mixed \\
M05 & Finance    & Household \\
M06 & Migration  & Individual \\
M07 & Health     & Individual \\
M08 & Food       & Household / individual \\
M09 & Hazards    & Household \\
M10 & Dwelling   & Household \\
M11 & Sanitation & Household \\
M12 & Utilities  & Household \\
M13 & Assets     & Household \\
M14 & Views      & Household \\
\bottomrule
\end{tabular}
\end{table}

The module-by-module rationale for each weight assignment is set out in the weight
rationale note.

\begin{refbox}
\rd{CAPI/Analysis/DB/doc/WEIGHT_RATIONALE.md}
\end{refbox}
